
\begin{proposition}
\label{XIII.1.14}
Sea $f\colon X \to S$ el morfismo estructural de un $S$-esquema, y sea
$F$ un haz de conjuntos o de grupos sobre~$X$ (\resp un haz de
ind-$\LL$-grupos, donde $\LL$ es un conjunto de números primos).
Supongamos que $F$ es localmente constante, que $(F,f)$ es
cohomológicamente propio en dimensión $\leqslant 0$ (\resp en dimensión
$\leqslant 1$) y que $f$ es localmente $0$-acíclico (\resp localmente
$1$-asférico para $\LL$) \textup{(SGA~4 XV~1.11)}. Entonces, para toda
especialización $\overline{s}_1 \to \overline{s}_2$ de puntos geométricos
de~$S$, el morfismo de especialización \textup{(SGA~4 VIII~7.1)}
$$
a_0 \colon (f_* F)_{\overline{s}_2} \to (f_*
F)_{\overline{s}_1}
$$
es un isomorfismo y, si $F$ es un haz de grupos, el morfismo
$$
a_1 \colon (\R^1 f_* F)_{\overline{s}_2} \to (\R^1 f_*
F)_{\overline{s}_1}
$$
es inyectivo (\resp los morfismos $a_0$ y $a_1$ son isomorfismos).
\end{proposition}

La
\marginpar{367}
demostración se obtiene reproduciendo palabra por palabra la de SGA~4
XVI~2.3, pero reemplazando la expresión \og propio\fg{} por la expresión
\og cohomológicamente propio\fg.

\ifthenelse{\boolean{orig}}
{}
{\enlargethispage{.5cm}}
